\documentclass{article}
\usepackage{listings}
\usepackage{xcolor}
\usepackage{enumitem}
\usepackage[margin=1in]{geometry}

\lstset{
  language=R,
  backgroundcolor=\color{gray!5},
  basicstyle=\ttfamily\small,
  keywordstyle=\color{black},
  commentstyle=\color{black},
  stringstyle=\color{black},
  showstringspaces=false,
  frame=single
}

% ---------- HEADER INFO ----------
\newcommand{\examtitle}{Midterm Exam}
\newcommand{\coursename}{040178-1: Applied Economics}
\newcommand{\semester}{Fall 2025}
\newcommand{\examdate}{Wednesday November 12th, 2025}
\newcommand{\duration}{60 minutes}

\begin{document}


\begin{center}
    {\Large \textbf{\examtitle}} \\[6pt]
    \coursename \\[4pt]
    \semester \\[4pt]
    \textbf{Date:} \examdate \quad | \quad \textbf{Duration:} \duration
\end{center}

\vspace{1em}
\noindent
\textbf{Instructions:}
\begin{itemize}
    \item This midterm exam is comprised of 25 multiple-choice questions.
    \item Each question features four possible answers. \textbf{Only one} answer is correct.
    \item Each question refers to the code chunk printed inmediately above it, unless otherwise stated.
    \item Assume that the code chunks are sequentially run, unless otherwise stated.
    \item The use of any external materials, including electronic devices, is sctrictly prohibited.
    \item In case you have a question, raise your hand and wait until the supervisor's permission to make it out loud.
    \item At the end of the exam, you should hand in:
    \begin{itemize}
      \item This question sheet.
      \item The ``Answers Form'' completed with your student ID number and your selected answers.
    \end{itemize}
\end{itemize}
\textbf{Grading Policy:}
\begin{itemize}
    \item This midterm exam represents 25\% of your final grade.
    \item Each correct answer adds +1 point.
    \item Each incorrect answer substracts 0.25 points.
    \item Each blank answer nor adds neither substracts points.
    \item Grading only depends on the ``Answers Form'' replies. You are allowed to sketch notes in this question sheets but these notes will not be taken into account for the final grade
    \item The use of any external materials, including electronic devices, is strictly prohibited.
\end{itemize}
\textbf{Good Luck!}
\newpage

%\section{R Objects}
\begin{lstlisting}
        n <- 1000 
        id <- 1:n 
        class(id)
        id[43]
        n[2]
\end{lstlisting}

\begin{enumerate}
  \item How many objects will there be in the ``Environment'' panel?
  \begin{enumerate}[label=\Alph*)]
    \item At least 1
    \item 2
    \item 1000
    \item At least 2
  \end{enumerate}
  \item What will the code line \texttt{id[43]} print in the ``Console'' panel?
  \begin{enumerate}[label=\Alph*)]
    \item \texttt{[1] 43}
    \item \texttt{[1] 42}
    \item \texttt{[1] 1 2 3 4 5 6}
    \item \texttt{NA}
  \end{enumerate}
  \item What will the code line \texttt{n[2]} print in the ``Console'' panel?
  \begin{enumerate}[label=\Alph*)]
    \item \texttt{[1] 1000}
    \item \texttt{[1] 1}
    \item \texttt{[1] 1 2 3 4 5 6}
    \item \texttt{NA}
  \end{enumerate}
\end{enumerate}

\newpage

\begin{lstlisting}
        rm(list = ls())
        n <- 100
        x <- rnorm(n, mean = -3, sd = 5)
        str(x)
        class(x)
        ls()
\end{lstlisting}

\begin{enumerate}[resume]
  \item How many objects will there be in the ``Environment'' panel?
  \begin{enumerate}[label=\Alph*)]
    \item 3
    \item At least 3
    \item 2
    \item At least 2
  \end{enumerate}
  \item What will the code line \texttt{str(x)} print in the ``Console'' panel?
  \begin{enumerate}[label=\Alph*)]
    \item \texttt{num [1:100]} \dots
    \item \texttt{num 100}
    \item \texttt{NA}
    \item \texttt{Error: object 'x' not found}
  \end{enumerate}
\end{enumerate}

\newpage

\begin{lstlisting}
  rm(list = ls())
  n <- 1000 
  brand   <- sample(c("BMW", "VW", "Toyota", "Renault", 
                      "Skoda", "Tesla"), n, replace = TRUE)
\end{lstlisting}

\begin{enumerate}[resume]
  \item How many arguments does function \texttt{sample()} take in this case?
  \begin{enumerate}[label=\Alph*)]
    \item Three
    \item Four
    \item Eight
    \item One 
  \end{enumerate}
  \item What is the role of argument \texttt{n} in function \texttt{sample()}?
  \begin{enumerate}[label=\Alph*)]
    \item Determine the number of observations that are drawn.
    \item Determine the population from which to sample.
    \item Determine the probability of selecting each element, at each draw.
    \item Determine whether sampling is with or without replacement.
  \end{enumerate}
  \item What is object \texttt{brand}?
  \begin{enumerate}[label=\Alph*)]
    \item A vector of with 1 element.
    \item A vector of 6 elements.
    \item A vector of 1000 elements.
    \item None of the above.
  \end{enumerate}
  \item What is the role of argument \texttt{replace = TRUE}?
  \begin{enumerate}[label=\Alph*)]
    \item Ensures each car brand (``BMW'', ``VW'', ``Toyota'' etc.) appears only once in object \texttt{brand}
    \item Enables for car brands to appear multiple times in object \texttt{brand}
    \item Ensures that object \texttt{brand} includes all car brands.
    \item None of the above statements are correct.
  \end{enumerate}
\end{enumerate}

\newpage

\begin{lstlisting}
  engine<- sample(c("Gasoline", "Diesel", "Hybrid", 
                    "Full Electric"), n, replace = TRUE,
                    prob = c(0.4, 0.35, 0.2, 0.05))
  table(engine)
\end{lstlisting}

\begin{enumerate}[resume]
  \item Which of the following statements is correct:
  \begin{enumerate}[label=\Alph*)]
    \item ``Gasoline'' is most likely the most frequent category in object \texttt{engine}
    \item ``Diesel'' is most likely the most frequent category in object \texttt{engine}
    \item ``Hybrid'' is most likely the least frequent category in object \texttt{engine}
    \item None of the above is correct.
  \end{enumerate}
  \item What is the most likely outcome of code line \texttt{table(engine)}?
    \begin{enumerate}[label=\Alph*)]
      \item 
      \begin{verbatim}
          Diesel Full Electric      Gasoline        Hybrid 
            84           125           718            73 
      \end{verbatim}

      \item 
      \begin{verbatim}
        Diesel Full Electric      Gasoline        Hybrid 
          346            48           413           193
      \end{verbatim}

      \item 
      \begin{verbatim}
          Diesel Full Electric      Gasoline        Hybrid 
            184           425           318            73 
      \end{verbatim}

      \item 
      \begin{verbatim}
          Diesel Full Electric      Gasoline        Hybrid 
            184           225            18           573 
      \end{verbatim}
    \end{enumerate}
\end{enumerate}

\newpage

\begin{lstlisting}
  index <- brand == "Tesla"
  engine[index] <- "Full Electric"
\end{lstlisting}

\begin{enumerate}[resume]
  \item What is object \texttt{index}?
  \begin{enumerate}[label=\Alph*)]
    \item A 6 element vector of logical \texttt{TRUE/FALSE} variables.
    \item A 1000 element vector of logical \texttt{TRUE/FALSE} variables.
    \item A 6 element vector of integer numbers.
    \item A 1000 element vector of integer numbers.
  \end{enumerate}
  \item What do the two code lines above do?
  \begin{enumerate}[label=\Alph*)]
    \item Change all the elements of vector \texttt{engine}
    \item Change some the elements of vector \texttt{engine}
    \item Change all the elements of vector \texttt{brand}
    \item Change some the elements of vector \texttt{brand}
  \end{enumerate}
\end{enumerate}

\newpage

\begin{lstlisting}
  age <- sample(0:15, n, replace = TRUE) 
  is_EV <- 1*(engine %in% c("Hybrid", "Full Electric"))
  price <- round(rnorm(n, mean = 40000 - 1500*age/10 + 10000*is_EV, sd = 8000), 0)
\end{lstlisting}

\begin{enumerate}[resume]
  \item Which of the following statements on \texttt{is\_EV} is correct?
  \begin{enumerate}[label=\Alph*)]
    \item Is a vector of logical \texttt{TRUE/FALSE} variables. The vector has around 250 elements.
    \item Is a vector of integer numbers. The vector has around 250 elements.
    \item It is a vector of logical \texttt{TRUE/FALSE} variables. The vector has 1000 elements.
    \item It is a vector of integer numbers. The vector has 1000 elements.
  \end{enumerate}
  \item How many arguments does \texttt{rnorm()} function take, in this case?
  \begin{enumerate}[label=\Alph*)]
    \item 4
    \item 5
    \item 9
    \item 3
  \end{enumerate}
  \item The object \texttt{price} will be comprised of\dots
  \begin{enumerate}[label=\Alph*)]
    \item logical \texttt{TRUE/FALSE} variables.
    \item String variables (words).
    \item Numbers with 2 decimals.
    \item Numbers with 0 decimals.
  \end{enumerate}
\end{enumerate}

\newpage

We next group the previous objects into a dataframe \texttt{df}. Assume variable price is in euros, and variable age is in years.

\begin{lstlisting}
  df <- data.frame(id, brand, engine, price, age, is_EV)
\end{lstlisting}

Let us start with the following two regression models \texttt{mlr\_0} and \texttt{mlr\_1}:
\begin{lstlisting}
  mlr_0 = lm( price ~ age , data = df )
  summary(mlr_0)
  #
  mlr_1 = lm( price ~ age + is_EV , data = df )
  summary(mlr_1)
\end{lstlisting}

\begin{enumerate}[resume]
  \item Based on the construction process of dataframe \texttt{df}, It is expected that the \texttt{mlr\_1} model\dots
  \begin{enumerate}[label=\Alph*)]
    \item ``age'' parameter estimate to be positive, ``is\_EV'' parameter estimate to be negative.
    \item ``age'' parameter estimate to be positive, ``is\_EV'' parameter estimate to be positive.
    \item ``age'' parameter estimate to be negative, ``is\_EV'' parameter estimate to be negative.
    \item ``age'' parameter estimate to be negative, ``is\_EV'' parameter estimate to be positive.
  \end{enumerate}
\end{enumerate}

\newpage
\begin{lstlisting}
  mlr_2 = lm( price ~ age*is_EV , data = df )
  summary(mlr_2)
\end{lstlisting}

\begin{enumerate}[resume]
  \item How many parameter estimates $\hat{\beta}_0$, $\hat{\beta}_1$, \dots, $\hat{\beta}_k$ will the multiple linear regression model \texttt{mlr\_2} have?
  \begin{enumerate}[label=\Alph*)]
    \item 7
    \item 6
    \item 2
    \item 4
  \end{enumerate}
  \item For a linear regression model $m$, $R^2_{\texttt{m}}$ assume denotes its R squared. How do you expect the R squares of models \texttt{mlr\_0}, \texttt{mlr\_1}, and \texttt{mlr\_2} to compare?
  \begin{enumerate}[label=\Alph*)]
    \item $R^2_{\texttt{mlr\_0}} \quad > \quad R^2_{\texttt{mlr\_1}} \quad = \quad R^2_{\texttt{mlr\_2}}$
    \item $R^2_{\texttt{mlr\_0}} \quad > \quad R^2_{\texttt{mlr\_1}} \quad \geq \quad R^2_{\texttt{mlr\_2}}$
    \item $R^2_{\texttt{mlr\_0}} \quad \leq \quad R^2_{\texttt{mlr\_1}} \quad \leq \quad R^2_{\texttt{mlr\_2}}$
    \item $R^2_{\texttt{mlr\_0}} \quad < \quad R^2_{\texttt{mlr\_1}} \quad = \quad R^2_{\texttt{mlr\_2}}$
  \end{enumerate}
\end{enumerate}

\newpage

\begin{lstlisting}
  for (i in c("BMW", "VW", "Toyota", "Renault", "Skoda", "Tesla")) {
        mlr_3 <- lm(price ~ age + is_EV, data = df[df$brand == i,])
        print(i)
        print(summary(mlr_3))
      }
\end{lstlisting}

\begin{enumerate}[resume]
  \item Which of the following statements is TRUE about this code?
  \begin{enumerate}[label=\Alph*)]
    \item The same dataset is used for each loop iteration.
    \item Each regression model is estimated using only observations of a specific brand.
    \item The \texttt{print(i)} statement is necessary for the code to work.
    \item The loop will generate 6 new objects, one for each car brand.
  \end{enumerate}
    \item What is the role of \texttt{print(summary(mlr\_3))}?

  \begin{enumerate}[label=\Alph*)]
    \item It prints only the regression coefficients.
    \item It provides detailed regression output.
    \item It provides detailed regression output, including R-squared and robust standard errors.
    \item It plots the residuals of the model.
  \end{enumerate}
\end{enumerate}

\newpage

\begin{lstlisting}
  df$log_price <- log(df$price)
  mlr_4 = lm(log_price ~ age , data = df )
  summary(mlr_4)
\end{lstlisting}
... prints the following parameter estimate in the RStudio Console:
\begin{verbatim}
  Coefficients:
                 Estimate   Std. Error    t value Pr(>|t|)    
  age           -0.003919     0.001313     -2.984  0.00292 ** 
\end{verbatim}

\begin{enumerate}[resume]
  \item Which of the following statements is correct?
  \begin{enumerate}[label=\Alph*)]
    \item An additional age unit is related to an expected price decrease of 0.39\%.
    \item An additional age unit is related to an expected price decrease of 0.39\%, holding other factors constant.
    \item An additional age unit is related to an expected price decrease of 39\%.
    \item An additional age unit is related to an expected price decrease of 39\%, holding other factors constant.
  \end{enumerate}
  \item The standard error \texttt{0.001313} estimation assumes homoskedasticity, this is \dots
  \begin{enumerate}[label=\Alph*)]
    \item All error terms $u_i$ must have the same conditional variance.
    \item All error terms $u_i$ may have the same conditional variance.
    \item All error terms $u_i$ must have different conditional variances.
    \item All error terms $u_i$ may have different conditional variances.
  \end{enumerate}
\end{enumerate}

\newpage

\begin{lstlisting}
  library(sandwich)  
  library(lmtest)   
  object <- vcovHC(mlr_2, type = "HC0")
  print(object)
\end{lstlisting}

\begin{enumerate}[resume]
  \item What do you expect the previous code to print in the Console?
  \begin{enumerate}[label=\Alph*)]
    \item The robust variances of the parameter estimates of model \texttt{mlr\_2}
    \item The robust standard errors of the parameter estimates of model \texttt{mlr\_2}
    \item The robust variance covariance matrix of the parameter estimates of model \texttt{mlr\_2}
    \item The robust variance covariance matrix of the parameter estimates of model \texttt{mlr\_2}, excluding intercept parameter.
  \end{enumerate}
  \item What are the dimensions of \texttt{object} ?
  \begin{enumerate}[label=\Alph*)]
    \item A $4 \times 4$ matrix.
    \item A $3 \times 3$ matrix.
    \item A four-element vector.
    \item A three-element vector.
  \end{enumerate}
\end{enumerate}

\end{document}
